% !TEX TS-program = pdflatexmk
\documentclass[12pt]{amsart}
\usepackage{multicol}
\usepackage{float}

%\usepackage[parfill]{parskip}    % Activate to begin paragraphs with an empty line rather than an indent

\usepackage[margin=1in]{geometry}


\usepackage{amsmath,amssymb,amsthm,latexsym,graphicx}
\usepackage[normalem]{ulem}
\usepackage{setspace} %used for doublespacing, etc.
\usepackage{hyperref}
\usepackage{cancel}
\usepackage[dvipsnames,usenames]{color}
\usepackage[all]{xy}
\usepackage{fancyhdr}
\pagestyle{fancy}
	\renewcommand{\headrulewidth}{0.5pt} % and the line
	\headsep=1cm
	
\DeclareGraphicsRule{.tif}{png}{.png}{`convert #1 `dirname #1`/`basename #1 .tif`.png}

%Some useful environments.
\newtheorem{theorem}{Theorem}
\newtheorem{corollary}[theorem]{Corollary}
\newtheorem{conjecture}[theorem]{Conjecture}
\newtheorem{lemma}[theorem]{Lemma}
\newtheorem{proposition}[theorem]{Proposition}
\newtheorem{definition}[theorem]{Definition}
\newtheorem{example}[theorem]{Example}
\newtheorem{axiom}{Axiom}
\theoremstyle{remark}
\newtheorem{remark}{Remark}
\newtheorem*{exercise}{Exercise}%[section]

%Some useful shortcuts for our favorite fields
\def\RR{\ensuremath{\mathbb R}} %Note, you can use these WITHOUT entering math mode
\def\NN{\ensuremath{\mathbb N}}
\def\ZZ{\ensuremath{\mathbb Z}}
\def\QQ{{\ensuremath\mathbb Q}}
\def\CC{\ensuremath{\mathbb C}}
\def\EE{{\ensuremath\mathbb E}}

%Some useful shortcuts for formatting lists
\newcommand{\bc}{\begin{center}}
\newcommand{\ec}{\end{center}}
\newcommand{\be}{\begin{enumerate}}
\newcommand{\ee}{\end{enumerate}}
\newcommand{\bi}{\begin{itemize}}
\newcommand{\ei}{\end{itemize}}

%Some useful shortcuts for formatting mathematical symbols
\newcommand{\ol}[1]{\overline{#1}}
\newcommand{\oimp}[1]{\overset{#1}{\iff}} %labeled iff symbol
\newcommand{\bv}[1]{\ensuremath{ \vec{\mathbf{#1}}} } %makes a vector.
\newcommand{\mc}[1]{\ensuremath{\mathcal{#1}}} %put something in caligraphic font
\newcommand{\mpg}[1]{\marginpar{ #1}} %to put comments in margins
\newcommand{\bsl}[1]{\texttt{\symbol{92}{\em #1}}} %for backslashes.
\newcommand{\normale}{\trianglelefteq}
\newcommand{\normal}{\triangleleft}

%Code for formatting the proofs a little nicer for submitted homework
\makeatletter
\renewenvironment{proof}[1][\proofname]{\par\doublespacing
  \pushQED{\qed}%
  \normalfont \topsep6\p@\@plus6\p@\relax
  \list{}{%
    \settowidth{\leftmargin}{\itshape\proofname:\hskip\labelsep}%
    \setlength{\labelwidth}{0pt}%
    \setlength{\itemindent}{-\leftmargin}%
  }%
  \item[\hskip\labelsep\itshape#1\@addpunct{:}]\ignorespaces
}{%
  \popQED\endlist\@endpefalse
  \singlespacing
}
\makeatother


%Commenting tools. You can ignore this stuff unless we're exchanging materials where I'm commenting in detail on how you are writing/using LaTeX.
\usepackage{soul}
\definecolor{highlight}{rgb}{1,0.6,0.6}
\sethlcolor{highlight}
\newcommand{\hlm}[1]{\colorbox{highlight}{$\displaystyle #1$}}
\newtheoremstyle{mycomment}{\topsep}{-0in}{\small \itshape \sffamily}{}{\small \itshape\sffamily}{:}{.5em}{}
\theoremstyle{mycomment}
\newtheorem*{acomment}{\color{BrickRed}{Comment}}
\newcommand{\com}[1]{{\color{OliveGreen}\begin{acomment}{#1} %#2 \color{black} 
\end{acomment}\noindent}}
%\newcommand{\com}[1]{{\color{BrickRed}{\\ Comment:}\color{OliveGreen}{#1} \\}}
\newcommand{\red}[1]{{\color{BrickRed} #1}}
\newcommand{\blue}[1]{{\color{MidnightBlue}#1}}
\newcommand{\green}[1]{{\color{OliveGreen}#1}}
\newcommand{\mwrong}[2]{\red{\cancel{#1}}\green{#2}}
\newcommand{\wrong}[2]{\red{\sout{#1}}\green{#2}}
\definecolor{OliveGreen}{rgb}{.3,.5,.2}
\definecolor{MidnightBlue}{rgb}{.3,.4,.6}
\newcommand{\pts}[1]{\hfill\blue{{#1}/5}}

\chead{MATH 631F}
\pagestyle{fancy}
%Modify these items:
\rhead{\emph{Stefano Fochesatto}}
\lhead{\emph{HW \#2 --- \today}}

\DeclareMathOperator{\lcm}{lcm}
\begin{document}

\thispagestyle{fancy}
\author{Coordinator: Coordinator's Name}





%Isomorphism Definition:
%phi : G \to H
%if phi is a homomorphism phi(xy) = phi(x)phi(y).
%
%phi is a bijection. 



\section*{\textbf{Section 1.6}}

%% Questionable Example
\begin{exercise}[\textbf{1.6.2}]If $\varphi: G \to H$ is an isomorphism, prove that $|\varphi(x)| = |x|$ for all $x \in G$. Deduce that any two 
  isomorphic groups have the same number of elements of order $n$ for each $n\in \ZZ^+$. Is the result true if $\varphi$ is only assumed to be a homomorphism.  
  \begin{proof} Let $\varphi: G \to H$ be an isomorphism and suppose $x \in G$ such that $|x| = \alpha$. By properties of homomorphisms,
    \begin{equation*}
        \varphi(x)^\alpha = \varphi(x^\alpha) = \varphi(e) = e. 
    \end{equation*}
    Suppose there is a $\beta < \alpha$ such that $\varphi(x)^{\beta} = e$. By properties of homomorphisms it would follow that, $x^\beta = e$ which is clearly a contradiction.\\
    Since $\varphi$ is a bijection, by the previous result if there are $m$ elements in $G$ of order $n$ there must be $m$ elements in $H$ of order $n$. This conclusion does not hold 
    if $\varphi$ is only a homomorphism, consider $\varphi: \ZZ_4 \to \ZZ_3$ such that $\varphi(x) = x \mod 3$. There are two elements in $\ZZ_4$ of order 5, and none in $\ZZ_3$ .
  \end{proof}
\end{exercise}
\vspace{.15in}

\begin{exercise}[\textbf{1.6.3}]If $\varphi: G \to H$ is an isomorphism, prove that $G$ is abelian if and only if $H$ is abelian. 
  if If $\varphi: G \to H$ is a homomorphism, what additional conditions on $\phi$ are sufficient to ensure that if $G$ is abelian, then so is $H$.\\
  \begin{proof} Let $\varphi: G \to H$ be an isomorphism, and suppose that $G$ is abelian. By definition for $a, b \in G$, $ab = ba$. Since $\varphi$ is 
    well defined, $\varphi(ab) = \varphi(ba)$ and by properties of homomorphisms, $\varphi(a)\varphi(b) = \varphi(b)\varphi(a)$. Since $\varphi$ is a surjection 
    every element in $H$ can be written in terms of $\varphi(x)$ for some $x \in G$. Thus $H$ is abelian.\\

    Now suppose $H$ is abelian. Since $\varphi$ is a surjection, consider $\varphi(a), \varphi(b) \in H$, and note $\varphi(a)\varphi(b) = \varphi(b)\varphi(a)$. By properties of homomorphisms, $\varphi(ab) = \varphi(ba)$
    and since $\varphi$ is a bijection we can apply $\varphi^{-1}$ to get, $ab = ba$. Thus $G$ is abelian.\\

    If $\varphi: G \to H$ is only a homomorphism, it is necessary for $\varphi$ to also be a surjection. This is needed when we define or note, $\varphi(a), \varphi(b) \in H$
    since every element in $H$ must be able to be written in terms of $\varphi(x)$ for $x \in G$. 
    This is awful yikers. 
  \end{proof} 
\end{exercise}
\vspace{.15in}

\begin{exercise}[\textbf{1.6.5}] Prove that the additive groups $\RR$ and $\QQ$ are not isomorphic.\\
  \begin{proof} Suppose for the sake of contradiction that $\phi: \RR \to \QQ$ is an isomorphism. By definition $\phi$ is also a bijection 
    from $\RR \to \QQ$ and therefore $|\RR| = |\QQ|$. Thus a contradiction. 
  \end{proof}
\end{exercise}
\vspace{.15in}


\begin{exercise}[\textbf{1.6.6}] Prove that the additive groups $\ZZ$ and $\QQ$ are not isomorphic.\\
  \begin{proof} Suppose for the sake of contradiction that $\phi: \QQ \to \ZZ$ is an isomorphism. By properties of isomorphisms we know that $\varphi(1) = 1$. Now consider $a \in \ZZ$. 
    by definition of $\ZZ$, $a = 1^a$. By substitution, and properties of homomorphisms $a = \varphi(1)^a = \varphi(1^a) = \varphi(a)$. Therefore we have shown that $\varphi(\ZZ) = \ZZ$ and thus 
    $\varphi$ cannot be an injection. 
  \end{proof}
\end{exercise}
\vspace{.15in}


\begin{exercise}[\textbf{1.6.14}] Let $G$ and $H$ be groups and let $\varphi: G \to H$ be a homomorphism. Define the kernel of $\varphi$ to be $\{g \in G | \varphi(g) = 1_H\}$.
  Prove that the kernel of $\varphi$ is a subgroup of $G$.
  \begin{proof} Let $G$ and $H$ be groups and let $\varphi: G \to H$ be a homomorphism. Suppose $g \in \ker(\varphi)$. By properties of isomorphisms, we know that $\varphi(1_G) = 1_H$. Consider that, 
    \begin{align*}
      \varphi(1_G) &= 1_H,\\
      \varphi(g^{-1}g) &= 1_H,\\
      \varphi(g^{-1})\varphi(g) &= 1_H,\\
      \varphi(g^{-1}) &= 1_H.
    \end{align*}
Thus $\ker(\varphi)$ is closed with respect to inverses.\\
Suppose $a,b \in \ker(\varphi)$. By properties of homomorphisms it follows that, 
\begin{align*}
  \varphi(a)\varphi(b) &= 1_H,\\
  \varphi(ab) &= 1_H.\\
\end{align*}
Therefore $\ker(\varphi)$ is closed. Thus $\ker(\varphi)$ is a subgroup of $G$.
\end{proof}


Prove that the kernel of $\varphi$ is injective if and only if the kernel of $\varphi$ is the identity subgroup of $G$.
\begin{proof}  Let $G$ and $H$ be groups and let $\varphi: G \to H$ be a homomorphism. Suppose that $ker(\varphi)$ is injective. It follows that $|ker(\varphi)| = 1$ since by definition 
  $|img(ker(\varphi))| = 1$. Note that,
  \begin{align*}
    \varphi(1_G) &= \varphi(1_G1_G),\\
    \varphi(1_G) &= \varphi(1_G)\varphi(1_G),\\
    \varphi(1_G)^{-1}\varphi(1_G) &= \varphi(1_G)^{-1}\varphi(1_G)\varphi(1_G),\\
    1_H &= \varphi(1_G).
  \end{align*}
  Therefore $1_G \in ker(\varphi)$. The other direction is trivial. 
\end{proof}
\end{exercise}
\vspace{.15in}



\begin{exercise}[\textbf{1.6.15}] Define the map $\pi: \RR^2 \to \RR$ by $\pi((x,y)) = x$. Prove that $\pi$ is a homomorphism and find the kernel of $\pi$. 
  \begin{proof} Consider the map $\pi: \RR^2 \to \RR$ by $\pi((x,y)) = x$. Suppose $(a, b),(c, d) \in \RR^2$. Note that 
    \begin{equation*}
      \pi((a, b) + (c, d)) = \pi((a + c, b + d)) = a + c = \pi((a, b)) + \pi((c, d)).
    \end{equation*}
    Note that $ker(\pi) = \{(a, b)\in \RR^2 | a = 0\}$.
  \end{proof}
\end{exercise}
\vspace{.15in}



\begin{exercise}[\textbf{1.6.17}] Let $G$ be any group. Prove that the map from $G$ to itself defined by $g \to g^{-1}$
  is a homomorphism if and only if $G$ is abelian.
  \begin{proof} Let $G$ be a group and $\varphi: G \to G$ such that $\varphi(g) = g^{-1}$. Suppose $\varphi$ is a homomorphism. Consider $a,b \in G$, by properties 
    of homomorphisms we know that, 
    \begin{equation*}
      ab = \varphi((ab)^{-1}) = \varphi(b^{-1}a^{-1}) = \varphi(b^{-1})\varphi(a^{-1}) = ba.  
    \end{equation*}
    Thus $G$ is abelian.\\

    Now suppose that $G$ is abelian. Consider $a,b \in G$,
    \begin{equation*}
      \varphi((ab)) = (ab)^{-1} = b^{-1}a^{-1} = a^{-1}b^{-1} = \varphi(a)\varphi(b). 
    \end{equation*}
    Thus $\varphi$ is a homomorphism.
  \end{proof}
\end{exercise}
\vspace{.15in}



\begin{exercise}[\textbf{1.6.18}] Let $G$ be any group. Prove that the map from $G$ to itself defined by $g \to g^2$
  is a homomorphism if and only if $G$ is abelian.
  \begin{proof} Let $G$ be a group and $\varphi: G \to G$ such that $\varphi(g) = g^{2}$. Suppose $\varphi$ is a homomorphism. 
    Consider $a,b \in G$, by properties 
    of homomorphisms we know that, 
    \begin{equation*}
      (ab)(ab) = (ab)^2 = \varphi(ab) = \varphi(a)\varphi(b) = a^2b^2.  
    \end{equation*}
    Applying $a^{-1}$ on the left and $b^{-1}$ on the right, we conclude that $ba = ab$. Thus $G$ is abelian. \\

    Now suppose that $G$ is abelian. Consider $a,b \in G$,
    \begin{equation*}
      \varphi((ab)) = (ab)^{2} = abab = a^2b^2 = \varphi(a)\varphi(b). 
    \end{equation*}
    Thus $\varphi$ is a homomorphism.
  \end{proof}
\end{exercise}
\vspace{.15in}



\begin{exercise}[\textbf{1.6.19}] Let $G = \{ z \in \CC | z^n = 1 \text{ for some } n \in \ZZ^+ \}$. Prove that for any fixed integer $k > 1$ 
  the map from $G$ to itself defined by $z \to z^k$ is a surjective homomorphism but is not an isomorphism.\\ 
  \begin{proof} 
  \end{proof}
\end{exercise}
\vspace{.15in}


\section*{\textbf{Section 1.7}}


\begin{exercise}[\textbf{1.7.3}] Show that the additive group $\RR$ acts on the $x, y$ plane $\RR x \RR$ by $r\dot(x, y) = (x + ry, y)$\\
  \begin{proof} Let $g_1, g_2 \in \RR$ and $(x, y) \in \RR x \RR$. Note that, 
     \begin{equation*}
      g_1 \dot (g_2 \dot (x, y)) = g_1 \dot  (x + g_2y, y) = (x + g_2y + g_1y, y) = (x + (g_2 + g_1)y, y) = (g_1 + g_2) \dot y.
     \end{equation*}
     Recall that in the additive group $\RR$, 0 is the identity. Note, 
     \begin{equation*}
      0\dot(x, y) = (x + 0y, y) = (x, y).
     \end{equation*}
     Thus the additive group $\RR$ acts on $\RR x \RR$.
  \end{proof}
\end{exercise}
\vspace{.15in}

\begin{exercise}[\textbf{1.7.4}] Let $G$ be a group acting on a set $A$ and fix some $a \in A$. 
  Show that the following sets are subgroups of $G$. 
  \begin{enumerate}
    \item[a.] The kernel of the action, 
    \begin{proof} Let $g, h$ be in the kernel of the action of $G$ and $x \in A$. Consider $gh$ acting on $x$, since 
      $g,h$ are in the kernel we get the following, 
      \begin{equation*}
        (gh)\cdot x = (g\cdot (h\cdot x)) = (g \cdot x) = x.
      \end{equation*} 
      Thus $gh$ is in the kernel of the action of $G$. Therefore the group is closed with respect to function composition. Note that,
      \begin{equation*}
        g^{-1}\cdot (x) = g^{-1}\cdot (g \cdot (x)) = (g^{-1}g) \cdot (x) = 1 \cdot x = x.
      \end{equation*}
      Therefore $g^{-1}$ is in the kernel of the action of $G$, therefore the group is closed with respect to inverses. Thus the 
      kernel of the action of $G$ is a subgroup of $G$. 
    \end{proof}
    \vspace{.15in}



    \item[b.] $ Stab_G(a) = \{g \in G | ga = a \}$. \\
    \begin{proof}Let $g, h \in Stab_G(a)$. Consider $gh$ acting on $a$, since 
      $g,h$ are in the stabilizer group of $a$ in $G$ we get, 
      \begin{equation*}
        (gh)\cdot a = (g\cdot (h\cdot a)) = (g \cdot a) = a.
      \end{equation*} 
      Thus $gh \in Stab_G(a)$ . Therefore the group is closed with respect to function composition. Note that,
      \begin{equation*}
        g^{-1}\cdot (a) = g^{-1}\cdot (g \cdot (a)) = (g^{-1}g) \cdot (a) = 1 \cdot a = a.
      \end{equation*}
      Therefore $g^{-1}\in Stab_G(a)$, therefore the group is closed with respect to inverses. Thus $Stab_G(a) \subseteq G$. 
      
    \end{proof}
  \end{enumerate}
\end{exercise}
\vspace{.15in}



\begin{exercise}[\textbf{1.7.6}] Prove that a group $G$ acts faithfully on a set $A$ if and only if the kernel of the action is the 
  set consisting only of the identity.\\
  \begin{proof} Suppose $G$, a group acting faithfully on the set $A$. Let $g$ be in the kernel of the action of $G$. By definition for 
    $a \in A$, we know that $g \cdot a = a$. Also note that $1 \cdot a = a$ and therefore $g$ and $1$ map to the same permutation representation.
    Since $G$ is a faithful action, by definition $\varphi: G \to S_A$ is injective, and therefore $g = 1$. Thus kernel of the action of $G$ 
    is only the identity. \\

    Suppose the kernel of the action is the set consisting only of the identity. For the sake of contradiction suppose that $G$ does not act faithfully on $A$, and therefore
    there exists $g, h \in G$ such that $g \neq h$ for all $a \in A$,  $g \cdot a = h \cdot a$. Note that 
    \begin{equation*}
      g^{-1}h \cdot (a) = g^{-1}\cdot (h \cdot (a)) = g^{-1}\cdot (g \cdot (a))  = (g^{-1}g) \cdot (a) = 1 \cdot a = a.
    \end{equation*}
    Note $g^{-1}h \neq 1$ since $g \neq h$, therefore $g^{-1}h$ in the kernel of the action is a contradiction.
  \end{proof}
\end{exercise}
\vspace{.15in}




\begin{exercise}[\textbf{1.7.13}] Find the kernel of the left regular action.\\
  \begin{proof} Let $G$ be a group, the kernel of it's left regular action is the set of all $g \in G$ such that 
    $ga = a$ for all $a \in G$. Clearly this set contains the identity element, and since $G$ is a group we know that the identity element is unique. 
    Thus the kernel of the left regular action is simply the identity.
  \end{proof}
\end{exercise}
\vspace{.15in}



%%%% Needs work 
\begin{exercise}[\textbf{1.7.14}] Let $G$ be a group and let $A = G$. Show that if $G$ is non-abelian then the maps defined 
  $g \cdot a = ag$ for all $g, a \in G$ do not satisfy the axioms of a left group action of $G$ on itself.\\
  \begin{proof} Let $G$ be non-abelian group and for the sake of contradiction suppose that the map $g \cdot a = ag$ for all $g, a \in G$ is a group action.
    By definition there exist $g, h \in G$ such that $gh \neq hg$. Note that 
    \begin{equation*}
      g \cdot (h \cdot a) = g \cdot (ah) = ahg 
    \end{equation*}
    By properties of group actions,
    \begin{equation*}
      g \cdot (h \cdot a) =   gh\cdot(a) = agh.
    \end{equation*}
    Therefore we get, $ahg = agh$ and by left cancellation of $a$ we get $hg = gh$ which is a contradiction. Thus $g \cdot a = ag$ does not satisfy the axioms of a left group action of $G$ on itself.
  \end{proof}
\end{exercise}
\vspace{.15in}

\begin{exercise}[\textbf{1.7.15}] Let $G$ be any group and let $A = G$. Show that the maps defined by $g \cdot a = ag^{-1}$ for all $g, a \in G$
  do satisfy the axioms of a left group action of $G$ on itself.\\
  \begin{proof} Let $G$ be any group and let $A = G$. Suppose the maps defined by $g \cdot a = ag^{-1}$ for all $g, a \in G$. Consider 
    \begin{equation*}
      (gh) \cdot a = a(gh)^{-1} = a(h)^{-1}(g)^{-1} = (g\cdot(h \cdot a)). 
    \end{equation*}
    Also note, 
    \begin{equation*}
      1 \cdot a = a (1)^{-1} = a
    \end{equation*}
    Therefore $g \cdot a = ag^{-1}$  defines a group action. 
  \end{proof}
\end{exercise}
\vspace{.15in}


\begin{exercise}[\textbf{1.7.16}] Let $G$ be any group and let $A = G$. Show that the maps defined by $g \cdot a = gag^{-1}$ for all 
  $g, a \in G$ do satisfy the axioms of a left group action of $G$ on itself. 
  \begin{proof} Let $G$ be any group and let $A = G$. Suppose the maps defined by $g \cdot a = ag^{-1}$ for all $g, a \in G$. Consider 
    \begin{equation*}
      (gh) \cdot a = (gh)a(gh)^{-1} = g(ha(h)^{-1})(g)^{-1} = (g\cdot(h \cdot a)). 
    \end{equation*}
    Also note, 
    \begin{equation*}
      1 \cdot a = (1) a (1)^{-1} = a
    \end{equation*}
    Therefore $g \cdot a = gag^{-1}$  defines a group action. 
  \end{proof}
\end{exercise}
\vspace{.15in}


\begin{exercise}[\textbf{1.7.17}] Let $G$ be a group and let $G$ act on itself by left conjugation, so each $g \in G$ maps $G$ to $G$
  by, 
  \begin{equation*}
    x \to gxg^{-1}
  \end{equation*}
  For fixed $g\in G$, prove that conjugation by $g$ is an isomorphism from $G$ into itself. Deduce that $x$ and $gxg^{-1}$ have the same order for 
  all $x$ in $G$ and that for any subset $A$ of $G$, $|A| = |gAg^{-1}|$.  
  \begin{proof} Suppose $g \in G$, and let $\varphi: G \to G$ be defined by $\varphi(x) =  gxg^{-1}$. Consider $a, b \in G$ and 
    note that, 
    \begin{equation*}
      \varphi(a)\varphi(b) = g(a)g^{-1}g(b)g^{-1} = g(a)(g^{-1}g)(b)g^{-1} = g(ab)g^{-1} = \varphi(ab). 
    \end{equation*}
    Therefore $\varphi$ is a homomorphism.\\
    Suppose $a, b \in G$ such that $\varphi(a) = \varphi(b)$. By left and right cancellation we get $a = b$ and therefore $\varphi$ is an injection. Since 
    $\varphi$ maps from $G$ to $G$ it is also a surjection. Therefore $\varphi$ is an isomorphism. \\
    By properties of isomorphisms, namely what was proved in exercise 1.6.2 we know that for all $x \in G$, $|\varphi(x)| = |x|$ and therefore $|x| = |gxg^{-1}|$.
    It also follows that for any $A \subseteq G$, $\varphi: A \to gAg^{-1}$ is an isomorphism, and therefore $|A| = |gAg^{-1}|$.
  \end{proof}
\end{exercise}
\vspace{.15in}


\begin{exercise}[\textbf{1.7.18}] Let $H$ be a group acting on a set $A$. Prove the the relation $\sim$ on $A$ defined by 
  \begin{equation*}
    a \sim b \text{ if and only if } a = h \cdot b \text{ for some $h \in H$}.
  \end{equation*}
  is an equivalence relation. 
  \begin{proof} Consider $a \in A$. Clearly $a \sim a$ since, $a = 1\cdot a$. Now suppose $b \in A$ such that $a \sim b$. By definition there exists some 
    $h \in H$ such that $a = h\cdot b$. Applying $h^{-1}$ we get, $ h^{-1}\cdot a = h^{-1}\cdot(h\cdot b) = (h^{-1}h)\cdot b = b$ and therefore $b \sim a$. 
    Suppose $a, b, c \in A$ such that $a \sim b$ and $b \sim c$. By definition we get that there exists some $g, h \in H$ such that $a = h \cdot b$ and $b = g \cdot c$. 
    By substitution we know that $a = g \cdot(h \cdot c) = (gh) \cdot c$ and therefore $a \sim c$. Thus $\sim$ is an equivalence relation.  
  \end{proof}
\end{exercise}
\vspace{.15in}


\begin{exercise}[\textbf{1.7.19}] Let $H$ be a subgroup of the finite group $G$ and let $H$ act on $G$ by left multiplication. 
  Let $x \in G$ and let $O$ be the orbit of $x$ under the action $H$. Prove the map 
  \begin{equation*}
    H \to O \text{ defined by } h \to h\cdot x
  \end{equation*}
  is a bijection. From this and the preceding exercise deduce Lagrange's Theorem. 
  \begin{proof} Let $\varphi: H \to O$ define the map $h \to hx$. Suppose $g, h \in H$ such that $\varphi(g) = \varphi(h)$.
    By definition we know that $gx = hx$ by right cancellation we conclude that $g = h$ and thus $\varphi$ is an injection. 
    Suppose $y \in O$. By definition there exists some $h \in H$ such that $y = hx$, and therefore $y = \varphi(h)$. Thus $\varphi$ is a surjection.\\

    Suppose a group $G$ and subgroup $H$. By exercise $1.7.18$ we know that the orbits under the action of $H$ partition $G$. Suppose there are $n$ orbits, so 
    it follows that, $|G| = \sum_{i = 1}^n |O|$. By the previous result since $H$ is a subgroup we know that all orbits, have cardinality $|H|$. Thus it follows that $|G| = n|H|$ and finally 
    $|H|\mid |G|$. 
    \end{proof}
\end{exercise}
\vspace{.15in}

\begin{exercise}[\textbf{1.7.20}] Show that the group of rigid motions of a tetrahedron is isomorphic to a subgroup of $S_4$
  \begin{proof} Suppose $G$ is the group of rigid motions of a tetrahedron. Let $A=[4]$ be the set of vertices. Now let $G$ act on $A$ and note that for each 
    $g \in G$ we get $\varphi_g: [4] \to [4]$, defined by $g \cdot a = \varphi_g(a)$ with $\varphi_g \in S_4$.
    Note that each $g \in G$ induces a different permutation on $a$ we know that the map from $G \to S_4$ defined $\sigma: g \to \varphi_g$ is an injection, and a homomorphism by properties of group actions. 
    Finally consider $\sigma(G)$ a subset $S_4$, which makes the map $G \to \sigma(G)$ defined $\sigma: g \to \varphi_g$ a surjection and therefore an isomorphism. 
    \end{proof}
\end{exercise}
\vspace{.15in}


\section*{\textbf{Section 2.1}}

\begin{exercise}[\textbf{2.1.3a}] Show that the $A = \{1, r^2, s, sr^2\}$ is a subgroup of the dihedral group $D_{2(4)}$.
    \begin{proof}[Solution] Since $D_{2(4)}$ is finite it is sufficient to show that $A$ is closed under composition for it to be a group. Recall the typical dihedral group presentation
      $D_{2(4)} = \{r, s|r^4 = s^2 = 1, rs = sr^{-1}\}$ and consider the following composition table, 
      
      \begin{figure}[H]
        \begin{tabular}{ c || c | c | c | c } 
          & 1 & $r^2$ & $s$ & $sr^2$\\
          \hline
          \hline
          1 & 1 & $r^2$ & $s$ & $sr^2$\\
          \hline
          $r^2$ & $r^2$ & 1 & $sr^2$ & $s$\\
          \hline
          $s$ & $s$ & $sr^2$ & 1 & $r^2$\\
          \hline
          $sr^2$ & $sr^2$ & $s$ & $r^2$& 1
        \end{tabular}.
      \end{figure}
      Clearly $A$ is closed under composition, therefore it is a subgroup of $D_{2(4)}$. 
    \end{proof}
\end{exercise}
\vspace{.15in}


\begin{exercise}[\textbf{2.1.4}] Give an explicit example of a group $G$ and an infinite subset $H$ of $G$ that is 
  closed under the group operation but is not a subgroup of $G$. 
    \begin{proof}[Solution] Any example must have the property that $H$ is not closed under inverses. 
      Consider $\RR^*$ under multiplication, and the set $\NN$. Clearly any product of two natural numbers is a natural number. The set also 
      contains the multiplicative identity, and is an infinite subset of $\RR$. However $\NN$ is not closed under inverses, and is therefore not a subgroup of $\RR^*$. 
  \end{proof}
\end{exercise}
\vspace{.15in}


\begin{exercise}[\textbf{2.1.5}] Prove that $G$ cannot have a subgroup $H$ with $|H| = n - 1$, where $|G| = n > 2$.
    \begin{proof} Let $G$ be a group with $|G| = n > 2$. Suppose $H$ is a subgroup of $G$ such that $|H| = n - 1$. 
      Since $G$ is finite and $H$ is a subgroup by Lagrange's Theorem we know that $|H| \mid |G|$, and therefore for some $i \in \ZZ$, $|G|=|H|i$. By substitution we get, 
      $n = (n - 1)i$ and therefore $i = n/(n - 1)$. Since $n>2$, there is no integer $i$ which satisfies $i = n/(n - 1)$, thus a contradiction.
    \end{proof}
\end{exercise}
\vspace{.15in}



\begin{exercise}[\textbf{2.1.6}] Let $G$ be an abelian group. Prove that $A = \{g\in G | |g| < \infty \}$ is a subgroup of $G$. 
  \begin{proof} Let $G$ be an abelian group. Suppose $A = \{g\in G | |g| < \infty \}$. Clearly $1 \in A$. Let $a, b \in A$, by definition 
    there exists some $m, n \in \NN$ such that $|a| = m$ and $|b| = n$. Consider $ab^{-1}$, since $G$ is abelian we know that, 
    \begin{equation*}
      (ab^{-1})^{mn} = (a^m)^n(b^n)^{m(-1)} = 1^n1^m = 1.
    \end{equation*}
    Thus $ab^{-1}$ must have finite order, and therefore $ab^{-1} \in A$. Thus $A$ is a subgroup of $G$.  
    \end{proof}
\end{exercise}
\vspace{.15in}



\begin{exercise}[\textbf{2.1.8}] Let $H$ and $K$ be subgroups of $G$. Prove that 
  $H \cup K$ is a subgroup if and only if either $H \subseteq K$ or $K \subseteq H$.
    \begin{proof} Let $H$ and $K$ be subgroups of $G$. Suppose that $H \cup K$ is a subgroup and without loss of generality let $H \not\subseteq K$. Therefore there exists some 
      $h \in H$ such that $h \not\in K$. Consider $k \in K$ and note that $hk \in H \cup K$ and therefore $hk \not\in K$ since $hk \in K$ implies $hk(k^{_1}) = h \in K$. Thus 
      $hk \in H$. It then follows that $(h^{-1}h)k = k \in H$ and therefore $K \subset H$.\\
      Without loss of generality suppose $K \subseteq H$. Therefore $H\cup K = H$, which is by definition a subgroup. 
    \end{proof}
\end{exercise}
\vspace{.15in}


\section*{\textbf{Section 2.2}}

\begin{exercise}[\textbf{2.2.5}] In each of the parts show that for the specified group $G$ and subgroup $A$ of $G$, 
  $C_G(A) = A$ and $N_G(A) = G$.
  \begin{enumerate}
    \item[a] $G = S_3$ and $A = \{1, (1 2 3), (1 3 2)\}$.
    \begin{proof}[Solution] 
    \end{proof} 

    \item[b] $G = D_{2(4)}$ and $A = \{1, s, r^2, sr\}$.
    \begin{proof}[Solution] 
    \end{proof} 

    \item[c] $G = D_{2(5)}$ and $A = \{1, r, r^2, r^3, r^4\}$.
    \begin{proof}[Solution]
    \end{proof}
  \end{enumerate}
\end{exercise}
\vspace{.15in}



\begin{exercise}[\textbf{2.2.6}] Let $H$ be a subgroup $G$. 
  \begin{enumerate}
    \item[a] Show that $H \leq N_G(H)$. Give an example to show that this is not necessarily true if $H$ is not a subgroup.\\
    \begin{proof} Let $H$ be a subgroup $G$. Suppose $h \in H$. Note that since $H$ is a subgroup 
      for all $a \in H$ we know that $hah^{-1} \in H$ and therefore by definition $h \in N_G(H)$ and thus $H \leq N_G(H)$.\\

      Let $G = S_3$, $H = \{1, (1 2 3) (1 3 2) (1 2)\}$. Note that $(1 2 3 )H(1 3 2) = \{1, (1 2 3) (1 3 2), (1 3)\} \neq H$ and therefore 
      $(1 2) \not\in N_G(H)$ so $H \not\leq G$.
    \end{proof} 

    \item[b] Show that $H \leq C_G(H)$ if and only if $H$ is abelian.\\
    \begin{proof} Let $H$ be a subgroup $G$. Suppose that $H \leq C_G(H)$. Consider $g, h\in H$ and note that since $H$ is a subgroup of $C_G(H)$ we 
      know that $ghg^{-1} = h$, therefore by right cancellation $gh = hg$. Thus $H$ is abelian.\\
      Now suppose $H$ us abelian. Consider $g, h \in H$ and note that since $gh = hg$ we know that $ghg^{-1} = h$. Therefore $g \in C_G(H)$ and thus $H \leq C_G(H)$.
    \end{proof} 
  \end{enumerate}
\end{exercise}
\vspace{.15in}




\begin{exercise}[\textbf{2.2.7}] Let $n \in \ZZ$ with $n \geq 3$. Prove the following:
  \begin{enumerate}
    \item[a]$Z(D_{2n}) = 1$ if $n$ is odd. 
    \begin{proof}
    \end{proof}

    \item[b]$Z(D_{2n}) = \{1, r^k\}$ if $n = 2k$.
    \begin{proof}
    \end{proof}
  \end{enumerate}
\end{exercise}
\vspace{.15in}



\begin{exercise}[\textbf{2.2.8}] Let $G = S_n$, fix an $i \in [n]$ and let $G_i = \{\sigma \in G | \sigma (i) = i \}$. Use group actions to prove that $G_i$ is 
  a subgroup of $G$. Find $|G_i|$.
    \begin{proof} Let $G = S_n$, fix an $i \in [n]$ and suppose $G_i = \{\sigma \in G | \sigma (i) = i$. Clearly $G_i$ is finite, therefore it suffices to show that 
      $G_i$ is nonempty and closed under function composition. Clearly $1 \in G_i$ by definition. Now consider $a, b \in G_i$, by the definition of our action $G$, 
      \begin{equation*}
        (ab)(i) =a(b(i)) = a(i) = i.  
      \end{equation*}
      Therefore $ab^{-1} \in G_i$ and $G_i$ is closed under function composition thus $G_i$ is a subgroup of $G$. 
    \end{proof}
\end{exercise}
\vspace{.15in}


\begin{exercise}[\textbf{2.2.9}] For any subgroup $H$ of $G$ and any nonempty subset $A$ of $G$ define $N_H(A)$ to be the set $\{h \in H | hAh^{-1} = A\}$. Show that 
  $N_H(A) = N_G(A)\cap H$ and deduce that $N_H(A)$ is a subgroup of $H$. 
    \begin{proof} Let $H$ be a subgroup of $G$ and for any nonempty subset $A$ of $G$ suppose $N_H(A)$ to be the set $\{h \in H | hAh^{-1} = A\}$. Since $H$ is 
      a subgroup of $G$ we know that $N_H(A) \leq N_G(A)$ and $N_H(A) \leq H$ then it follows that $N_H(A) \neq N_G(A)\cap H$ . Let $g \in N_G(A) \cap H$, by definition $gAg^{-1} = A$ and $g \in H$. Thus $g \in N_H(A)$. Therefore $N_H(A) = N_G(A)\cap H$.\\

      Clearly $N_H(A)$ is nonempty since $1A1 = A$. Also $N_H(A)$ is closed with respect to inverses, for $h \in N_H(A)$ we know that $hAh^{-1} = A$. By left and right cancellation we get, $A= h^{-1}Ah$ therefore $h^{-1} \in N_H(A)$.
      Finally $N_H(A)$ is closed with respect to composition, let $a, b \in  N_H(A)$ and consider that $abA(ab)^{-1} = a(bAb^{-1})a^{-1} = a(A)a^{-1} = A$ therefore $ab \in N_H(A)$. Thus $N_H(A)$ is a subgroup. 
    \end{proof}
\end{exercise}
\vspace{.15in}

\begin{exercise}[\textbf{2.2.10}] Let $H$ be a subgroup of order 2 in $G$. Show that $N_G(H) = C_G(H)$. Deduce that if $N_G(H) = G$ then $H \leq Z(G)$.
    \begin{proof} 

    \end{proof}
\end{exercise}
\vspace{.15in}

\begin{exercise}[\textbf{2.2.11}] Prove that $Z(G) \leq N_G(A)$ for any subset $A$ of $G$.
  \begin{proof} Let $A$ be a subset of $G$. Suppose $g \in Z(G)$. By definition $gx = xg$ for all $x \in G$. By right cancellation it follows that 
    $gxg^{-1} = x$ for all $x\in G$. Since $A$ is a subset of $G$, it follows that $gag^{-1} = a$ for all $a \in A$ and therefore $g\in  N_G(A)$. Thus
    $Z(G) \leq N_G(A)$. 
  \end{proof}
\end{exercise}
\vspace{.15in}


\end{document} 



